% Generated draft: requires mathematical and source review.
\documentclass[11pt]{article}
\usepackage[margin=25mm]{geometry}
\usepackage{fontspec}
\setmainfont{Latin Modern Roman}
\usepackage{amsmath,amssymb,mathtools,unicode-math}
\setmathfont{Latin Modern Math}
\usepackage{enumitem,needspace,xurl,hyperref}
\setlength{\parindent}{0pt}
\setlength{\parskip}{5pt}
\setlength{\emergencystretch}{3em}
\newcommand{\N}{\mathbb{N}}
\newcommand{\Z}{\mathbb{Z}}
\newcommand{\Q}{\mathbb{Q}}
\newcommand{\R}{\mathbb{R}}
\newcommand{\C}{\mathbb{C}}
\newcommand{\F}{\mathbb{F}}
\newcommand{\A}{\mathbb{A}}
\newcommand{\PP}{\mathbb{P}}
\newcommand{\E}{\mathbb{E}}
\newcommand{\eps}{\varepsilon}
\newcommand{\dd}{\,\mathrm{d}}
\begin{document}
\section*{UnsolvedMath ID: 30005356}
\textbf{Source problem code:} catalogue entry
\subsection{Definitions and mathematical statement}
The supplied record asks about a \emph{generic multiplicative field} $(K,\sigma)$, but does not specify the axioms of genericity, the characteristic of $K$, the domain of the endomorphisms, or the parameters permitted in definitions. Consequently, the following is a conditional formulation under the multiplicative-group interpretation; these choices require source verification.

Let $K$ be a field, write $K^\times=K\setminus\{0\}$, and suppose that $\sigma:K^\times\to K^\times$ is a group endomorphism. In the ring $\operatorname{End}(K^\times)$, addition is pointwise multiplication and multiplication is composition:
\[
(f+g)(x)=f(x)g(x),\qquad (fg)(x)=f(g(x)).
\]
The multiplicative identity of this ring is $\operatorname{id}_{K^\times}$, and the integer $n$ acts as $x\mapsto x^n$. Thus the subring generated by $\sigma$ and the integer scalar endomorphisms is
\[
\mathbb Z[\sigma]
=\left\{x\longmapsto\prod_{i=0}^{m}\bigl(\sigma^i(x)\bigr)^{n_i}:
 m\in\mathbb Z_{\geq0},\ n_0,\ldots,n_m\in\mathbb Z\right\},
\qquad \sigma^0=\operatorname{id}_{K^\times}.
\]
Here $\sigma^i$ denotes iteration, and negative exponents are meaningful because the domain is $K^\times$.

Fix an intended generic structure $(K,\sigma)$, its intended first-order language, and a set $A\subseteq K$ of permitted parameters. A map is $A$-definable if its graph is defined in that structure by a formula using finitely many parameters from $A$. The question under these conventions is whether every $A$-definable group endomorphism $f:K^\times\to K^\times$ satisfies
\[
\exists m\in\mathbb Z_{\geq0}\ \exists(n_0,\ldots,n_m)\in\mathbb Z^{m+1}\ \forall x\in K^\times,
\qquad f(x)=\prod_{i=0}^{m}\bigl(\sigma^i(x)\bigr)^{n_i}.
\]
The coefficients and degree may depend on $f$. This question is to be posed for the structures intended by the source. The supplied record does not determine that class, the language, or the allowed parameter convention, and does not establish that $K^\times$ is the intended domain. These unresolved qualifications prevent an unconditional formalization of the extracted question.
\subsection{Short English statement}
Under the multiplicative-group interpretation, the question asks whether every definable endomorphism is a finite product of integer powers of iterates of $\sigma$. The meaning of genericity, the permitted parameters, and the endomorphism domain remain unspecified. The supplied dated curation labels this question open; that status is not independently verified here.
\subsection{Sources}
\begin{itemize}
\item Supplied local UnsolvedMath registry record, permanent ID $30005356$, title \emph{Definable Endomorphisms of Generic Multiplicative Fields}, source problem code \texttt{catalogue entry}, in the Oberwolfach Reports Open Problems collection. The record gives proposed year $2023$ and no named proposer. Its original passage attributes the question to ``the author'' without identifying that person.
\item \emph{Model Theory: Combinatorics, Groups, Valued Fields and Neostability} (2023), Oberwolfach Reports, pp.~91--149, DOI 10.4171/owr/2023/2. Supplied source citation and literature reference: \url{https://doi.org/10.4171/owr/2023/2}. No precise page or passage locator for this question is supplied.
\item Mutchnik, \emph{An NSOP1 Theory without the Existence Axiom}. Supplied literature reference: \url{https://arxiv.org/abs/2407.13082}.
\item d'Elbee, \emph{Generic Multiplicative Endomorphism of a Field}. Supplied literature reference: \url{https://doi.org/10.1016/j.apal.2025.103554}.
\item Literature assessment dated exactly \texttt{2026-08-22}: the supplied curation states that Mutchnik disproves universal forking existence in NSOP1 theories, making the first sentence of the original passage outdated, while the bundled definable-endomorphism question remains unresolved. This is supplied dated curation, not an independently verified literature assessment. The record itself requests independent verification of time-sensitive status and cited claims.
\item The registry labels the statement \texttt{corrected\_verified}, explaining that only formula delimiters, notation, grammar, and typography were repaired and that mathematical content and scope were preserved. This supplied label does not resolve the formulation gaps or constitute mathematical review of this draft.
\end{itemize}
\end{document}
